\section{One selector for spatial and semifinite perfection}\label{sec:iionly}
The critical selector can choose the same witness in spatial and semifinite
models. The construction uses two edge families on one scaffold: the
selector is critical, while the four-letter computation tree is subcritical.
The key point is that a nonzero selected terminal sector has positive
intersection with the computation root. This supplies a finite witness
without using a density eigenvector in a diffuse host.

\begin{lemma}[Several edge families on a common chain]\label{lem:multifamily}
Fix $n\ge1$ and finitely many edge families $X_i^{(r)}$. Each family has
its own Gram equations, initial defect and range equation, and its own
within-family endpoint equations with typed centralizer coefficients.
There are no prescribed cross-family Gram equations. If all edges have
modular ratio $2^{-n}$, one finite rational Boolean game compiles these
systems on a common nonzero source corner. It extracts them from every
perfect normal commuting strategy, and every faithful normal realization
extends simultaneously, preserving host type under finite amplification
and finite entropy for spatial realizations.
\end{lemma}
\begin{proof}
Use the corners, activities, bound axes, and common chain
$v=\btrans_n\cdots\btrans_1$ of Lemma~\ref{lem:chain}. Its source mass is
$\mathfrak c_n^{-1}$. For each edge, independently introduce a supported
centralizer contraction $a_i^{(r)}=e_na_i^{(r)}e_0$. A family with initial
defect $D_r$ and range $R_r$ is submitted as
\[
 a_i^{(r)*}a_j^{(r)}=\delta_{ij}(e_0-(D_r)_0),\qquad
 \sum_i a_i^{(r)}a_i^{(r)*}=(R_r)_n.
\]
Only that family's own equations are submitted. Transport the union of
all labels appearing at left endpoints or in range equations, with one
separate graph refinement per encoding and bridge. The common original
signs force $L_nv=vL_0$ for every label in this union. Decode
$X_i^{(r)}=v^*a_i^{(r)}$. Multiplication by $v,v^*$ proves the Gram, range
and endpoint translations in both directions; for example
$L_na_i^{(r)}=a_j^{(r)}B_0$ becomes
$L_0X_i^{(r)}=X_j^{(r)}B_0$. The centralizer condition and the modular
eigenvalue of $v$ give ratio $2^{-n}$ for every extracted edge.
All equations concern the same faithful corner and state.

Conversely use $\operatorname{diag}(\mu_n)\otimes\varphi_0$ and
$a_i^{(r)}=|e_n\rangle\langle e_0|\otimes X_i^{(r)}$ for every family.
Every lifted contraction lies in the finite centralizer and has its
own two-unitary encoding. Choose the root encodings once and copy them
through the identity corner maps. The resulting block labels commute
with their bridge's canonical signs. Actual joint products give all
separate refinements, followed by the common qutrit and row extensions.
There is one state and one assignment for all families. The construction
adds only finite factors and imposes no mixed orthogonality.
\end{proof}

\subsection{The coupled root presentation}
Use four edges $A_{ab}$ satisfying \eqref{eq:fockrelations} at modular
ratio $1/16$, with root $P_\circ$, and sixteen different edges $S_i$ with
\[
 S_i^*S_j=\delta_{ij}(I-P),\quad \sum_{i=1}^{16}S_iS_i^*=I-Q,
 \quad \sigma_t^\varphi(S_i)=16^{-it}S_i.
\]
For a positive centralizer contraction $C_{\rm sel}$ impose
\[
 C_{\rm sel}Q=Q,\quad C_{\rm sel}S_i=\tfrac12S_iC_{\rm sel},
 \quad[C_{\rm sel},P]=0,\quad[P,A_{ab}]=[C_{\rm sel},A_{ab}]=0.
\]
Lemma~\ref{lem:selector} above uses modular ratio $1/16$ and
depth ratio $1/2$. Retain the coordinate and boundary equations
\eqref{eq:U_x}--\eqref{eq:P_y} and the counter equations
\eqref{eq:C}, with $[C_{\rm ctr},C_{\rm sel}]=0$. Replace the clock by
\begin{equation}\label{eq:equalityreader}
 (2C_{\rm ctr}-C_{\rm sel})P_{\rm ret}=0,\qquad
 T_{\rm inv}(2C_{\rm ctr}-C_{\rm sel})=
 \tfrac18C_{\rm ctr}C_{\rm sel}(I-P_{\rm ret}).
\end{equation}
Here $T_{\rm inv}$ is a general contraction. The tableau has the same
initial and update equations, but its halt prohibition is
$P_{\mathtt H}P=0$. Its machine reads the first positive pulse ${d_{\rm clk}}$, sets
$n={d_{\rm clk}}-1$, and tests $\Rdec(e,n,0),\Rdec(e,n,1),\ldots$, halting at the
first false answer. No prime test is used.

Lemma~\ref{lem:multifamily} with $n=4$ compiles the two separate Gram
systems. The transported typed list is exactly
\[
 P_\circ,U_x,U_y,P_x,P_y,C_{\rm ctr},P,C_{\rm sel},Q.
\]
The return projection, its inverse and the tableau projections are
source-only. In particular no relation $A_{ab}^*S_i=0$ is imposed.
Denote this finite game by $F_e$.

\subsection{The terminal sector meets the root}
\begin{lemma}\label{lem:terminalroot}
Let orthogonal isometries $A_1,\ldots,A_r$ have modular ratio
$\nu$ on a faithful normal support. If a centralizer projection $p$
commutes with them and $r\nu<1$, then
\[
 \varphi\left(p\left(I-\sum_i A_iA_i^*\right)\right)
 =(1-r\nu)\varphi(p).
\]
\end{lemma}
\begin{proof}
Commutation with a self-adjoint projection also gives commutation with
the adjoints. KMS gives
$\varphi(A_ipA_i^*)=\nu\varphi(pA_i^*A_i)=\nu\varphi(p)$.
Subtract the sum of these $r$ terms from $\varphi(p)$.
\end{proof}

In a semifinite realization, Lemma~\ref{lem:selector} with $d=16$ gives
some nonzero $p_n=P1_{\{2^{-n}\}}(C_{\rm sel})$. It commutes with all
$A_{ab}$ and their adjoints. Hence
\begin{equation}\label{eq:terminalroot}
 \varphi(p_nP_\circ)=\tfrac34\varphi(p_n)>0.
\end{equation}
Choose $0\ne\xi\in p_nP_\circ\HH$. Its word descendants remain in
$p_n\HH$, are nonzero, and are orthogonal at every fixed word length.
The coordinate inductions identify every finite square on these vectors.

For self-adjoint $X,Y$, with $Y$ injective, the two equations
$(2X-Y)R=0$ and $T(2X-Y)=\kappa XY(I-R)$, $\kappa\ne0$, force
$R=1_{\{0\}}(2X-Y)$. Indeed $R\le K=1_{\{0\}}(2X-Y)$.
For $\zeta\in K\HH$ put $\eta=(I-R)\zeta\in K\HH$. The second equation
gives $XY\eta=0$, and $Y\eta=2X\eta$ gives $2X^2\eta=0$.
Self-adjointness gives $X\eta=0$, then $Y\eta=0$, then $\eta=0$.
Apply this to \eqref{eq:equalityreader}, using injectivity of
$C_{\rm sel}$ from the selector lemma. On the chosen descendants,
$C_{\rm ctr}=2^{-y}$ and $C_{\rm sel}=2^{-n}$. Thus the pulse is
exactly $y=n+1$, including no time-zero pulse.

\begin{lemma}[Definite computation on a terminal sector]\label{lem:terminaltableau}
Suppose $0\ne\xi\in P_\circ\HH$, every $A_w\xi$ belongs to $P\HH$,
and $P_{\rm ret}$ has the specified classical scalar on each $A_w\xi$.
With the coordinate, initial and update equations above, every symbol
projection acts on these descendants as the genuine classical symbol.
If the presentation imposes $P_{\mathtt H}P=0$, that classical
computation cannot halt.
\end{lemma}
\begin{proof}
At each fixed depth, the coordinate recursions give cyclic shifts and
boundary masks on the word descendants. Row zero follows from
\eqref{eq:initial} without its global halt equation. Induct on time at
all spatial coordinates simultaneously. Each factor in the displayed
ordered update product acts as its classical zero or one on the
induction vector: left and right factors refer to the preceding row,
and the pulse is scalar by hypothesis. Exactly the prescribed successor
survives. Multiplication by $U_y$ gives the next-row assertion. The
last-row mask removes wraparound; radius one and blank quiescence make
the right ghost blank throughout the required prefix. This proves the
half-line computation on every sufficiently large square. If a halt
occurred, one nonzero descendant $\eta\in P\HH$ would satisfy
$P_{\mathtt H}\eta=\eta$, contradicting
$P_{\mathtt H}\eta=P_{\mathtt H}P\eta=0$.
The conclusion uses the restricted prohibition, not a claim that
$P_{\mathtt H}=0$ on the entire space.
\end{proof}

\begin{theorem}[A shared family]\label{thm:sharedfamily}
The family $F_e$ satisfies
\[
 F_e\in\LE\iff F_e\in\LI\iff F_e\in\LU
 \iff\exists n\,\forall t\,\Rdec(e,n,t),\qquad F_e\notin\LJ.
\]
\end{theorem}
\begin{proof}
Semifinite soundness follows from the selected sector, the exact pulse
and Lemma~\ref{lem:terminaltableau}. A false check at that selected $n$
would halt, so $n$ is a witness.

For completeness let $n$ be a witness. On the finite sixteen-letter
tree of words $v$ of length at most $n$, use prefix maps $S_i$, zero
on leaves, with $P$ the leaves, $Q$ the root and
\[
 h_{\rm sel}|v\rangle=(n+1)^{-1}16^{-|v|}|v\rangle,
 \qquad C_{\rm sel}|v\rangle=2^{-|v|}|v\rangle.
\]
Tensor it with the four-letter Fock density $h_F$ of
\eqref{eq:fockdensity}. At selector depth $\ell$, put
$P_{\rm ret}=1_{\{y=\ell+1\}}$ and
\[
 T_{\rm inv}(y,\ell)=
 \begin{cases}[8(2^{\ell+1}-2^y)]^{-1},&y\ne\ell+1,\\
 0,&y=\ell+1.\end{cases}
\]
The nonzero denominator before the factor eight is an integer. The norm
is exactly $1/8$, attained at $\ell=0,y=0$. Negative values occur.
Use the genuine classical tableau for input $\ell$ on each such fibre;
all symbol projections are its diagonal indicator projections. Coordinate
permutations preserve Fock depth, and every other centralizer label is
diagonal. Both edge families have ratio $1/16$. Computations at depths
$\ell<n$ may halt; the terminal equation excludes only depth $n$.
This gives every label and equation simultaneously, also at $n=0$,
where $S_i=0$, $P=Q=C_{\rm sel}=I$.

The selector entropy is $\log_2(n+1)+2n$, so the root product entropy is
\[
 \log_2(n+1)+2n+\log_2(4/3)+4/3<\infty.
\]
The compiler adds finite entropy. The inclusions
$\LE\subseteq\LI\subseteq\LU$ finish the equivalence. In a finite
host the four extracted isometries would have orthogonal range traces
summing to four under a normalized faithful trace, but their sum is at
most the identity. Thus every instance is outside $\LJ$.
\end{proof}
